% Generated draft: requires mathematical and source review.
\documentclass[11pt]{article}
\usepackage[margin=25mm]{geometry}
\usepackage{fontspec}
\setmainfont{Latin Modern Roman}
\usepackage{amsmath,amssymb,mathtools,unicode-math}
\setmathfont{Latin Modern Math}
\usepackage{enumitem,needspace,xurl,hyperref}
\setlength{\parindent}{0pt}
\setlength{\parskip}{5pt}
\setlength{\emergencystretch}{3em}
\newcommand{\N}{\mathbb{N}}
\newcommand{\Z}{\mathbb{Z}}
\newcommand{\Q}{\mathbb{Q}}
\newcommand{\R}{\mathbb{R}}
\newcommand{\C}{\mathbb{C}}
\newcommand{\F}{\mathbb{F}}
\newcommand{\A}{\mathbb{A}}
\newcommand{\PP}{\mathbb{P}}
\newcommand{\E}{\mathbb{E}}
\newcommand{\eps}{\varepsilon}
\newcommand{\dd}{\,\mathrm{d}}
\begin{document}
\section*{UnsolvedMath ID: 800978}
\textbf{Source problem code:} catalogue entry
\subsection{Definitions and mathematical statement}
Let \(M\) be a smooth surface of genus at least \(2\). Under the usual smooth, strongly convex convention, a Finsler metric is a continuous function \(F:TM\to[0,\infty)\), smooth on \(TM\setminus\{0\}\), positive on nonzero tangent vectors, and satisfying \(F(x,\lambda y)=\lambda F(x,y)\) for every \(\lambda>0\), whose fundamental tensor
\[
g_{ij}(x,y)=\frac12\frac{\partial^2(F^2)}{\partial y^i\partial y^j}(x,y)
\]
is positive definite for every \((x,y)\in TM\setminus\{0\}\). Define its Cartan tensor and Landsberg tensor by
\[
C_{ijk}=\frac12\frac{\partial g_{ij}}{\partial y^k},\qquad L_{ijk}=C_{ijk\mid s}y^s,
\]
where \(\mid\) denotes horizontal covariant differentiation for the Chern connection and repeated indices are summed. The Landsberg condition is \(L\equiv0\). The metric is Riemannian if \(F(x,y)^2=h_x(y,y)\) for a Riemannian metric \(h\); under these conventions, being non-Riemannian is equivalent to \(C\not\equiv0\).

The supplied Question 59 asks whether there exist such a surface \(M\) and a non-Riemannian Finsler metric \(F\) satisfying
\[
\operatorname{genus}(M)\ge2,\qquad L\equiv0,\qquad C\not\equiv0.
\]
The supplied excerpt does not specify whether the surface must be closed, compact, orientable, or without boundary, nor does it specify reversibility or the precise regularity convention for Finsler metrics. These qualifications remain unresolved; the definitions above state the usual smooth, strongly convex interpretation without resolving the intended surface class or imposing reversibility.

The preceding supplied paragraph concerns an unavailable previous question: it attributes to [Mat-Top1998] a construction of an integral of one geodesic flow using a second metric, and states that in the Riemannian case the integral is quadratic in velocities and an affirmative answer follows from [Kolokoltsov1983], with [Mat-Top2000] also cited. Its hypotheses and the target of that affirmative answer are not supplied, so this contextual assertion is not an answer to the existence question formulated here.
\subsection{Short English statement}
Does a surface of genus at least \(2\) admit a non-Riemannian Finsler metric with vanishing Landsberg tensor? The intended surface class and some metric conventions are unspecified. The supplied registry labels the question open; this draft does not independently verify that status.
\subsection{Sources}
\begin{itemize}
\item Supplied local UnsolvedMath registry record: permanent ID \(800978\), exact source problem code \texttt{catalogue entry}, title ``Open Problems and Questions about Geodesics --- Question 59.'' Attribution supplied by the registry: Burns, Matveev; proposed year 2019; supplied status: open.
\item Burns, Matveev, \emph{Open Problems and Questions about Geodesics} (2019), Question 59, as identified in the supplied source record: \url{https://arxiv.org/abs/1308.5417}. The record reports access on 2026-10-16 and extraction by \texttt{arxiv\_latex}; these are supplied provenance statements, not an independently verified source visit.
\item Contextual reference labels supplied in the excerpt: [Mat-Top1998], [Kolokoltsov1983], and [Mat-Top2000]. Full bibliographic entries and their precise hypotheses were not supplied.
\end{itemize}
\end{document}
